\fchapter{Ej 7 Pág 78}

\fsection{\boldit{\textcolor{waveBlue2}{Investiga.}}} 
\subsection*{En Suecia existe una tradición cultural llamada FIKA. Realiza una búsqueda en internet sobre este fenómeno y elabora una breve reseña sobre las aportaciones que contribuyen a la consecución del ODS 8.}

\begin{solution}
La Fika es una pausa social fundamental en la cultura sueca, que consiste en tomar un café (o infusión) 
acompañado de un dulce, en compañía de compañeros de trabajo, amigos o familiares. Va más allá de un simple 
descanso; es un ritual integrado en la jornada laboral y social que promueve el bienestar.

\textbf{Aportaciones clave a la consecución del ODS 8 (Trabajo Decente y Crecimiento Económico):}
\begin{enumerate}
	\item \textbf{Fomenta un entorno de trabajo saludable y productivo: Al institucionalizar pausas regulares y desconexión, la Fika reduce el estrés, previene el agotamiento (burnout) y mejora la concentración. Esto aumenta la productividad y la creatividad, contribuyendo a un crecimiento económico sostenible.}
	\item \textbf{Refuerza la cohesión de equipo y la comunicación informal: Esta pausa compartida nivela las jerarquías, facilita la resolución informal de problemas y fortalece los lazos entre colegas. Un buen clima laboral mejora la retención del talento y reduce la rotación.}
	\item \textbf{Promueve el equilibrio entre vida laboral y personal (conciliación): La Fika marca una separación psicológica dentro de la jornada. Su carácter social y de ocio dentro del horario laboral subraya la importancia del bienestar personal, alineándose con la meta de un "trabajo decente".}
	\item \textbf{Sostiene la economía local: La tradición impulsa el consumo regular en cafeterías locales y panaderías, apoyando a pequeños negocios y creando empleo en el sector servicios.}
\end{enumerate}
\textbf{En resumen,} la Fika es un ejemplo práctico de cómo una práctica cultural aparentemente simple puede tener un impacto profundo en la calidad del trabajo, 
la productividad económica y el bienestar social, ejemplificando el espíritu del ODS 8.

\end{solution}
